% Français 9115, batch 17 — Gallica views 321–340.
% Pass: Opus 5.5 (claude-opus-5-5), 2026-09-22 — first pass, unchecked against the views by a human.
%
% What the twenty views hold. One piece, in a fair copy: a memoir in
% French headed, at view 320 (the last view of batch 16, not transcribed
% here), « Remarques sur le mémoire d'Euler: Investigatio motuum quibus
% laminæ et virgæ elasticæ contremiscunt. in act acad Petrop ann 1779 P 1
% p. 103 et seq », which is a new and much more finished redaction of the
% critique of Euler's § 47 that filled views 208–299. It does not continue
% the draft memoir of batch 15 (whose last page, view 299, stopped
% mid-sentence on « les deux »): view 320 opens a new series of pages,
% (1), and view 322 is (2). The hand is a regular, calligraphic fair hand,
% almost without correction — very unlike the heavily corrected hand of
% views 218–299; this pass does not say whether it is hers or a
% copyist's, although the text writes « afin d'être autorisée » in the
% first person (view 332). Sentences of the draft recur here nearly word
% for word (« J'ai voulu entrer dans tous ces détails, àfin d'être
% autorisée a conclure qu'Euler a été induit en erreur », view 332, cf.
% view 240), and the arguments of the draft (the second solution
% incompatible with sin ω = 0; nodes at the aliquot parts; the rod free
% at both ends) come back in order and in cleaner form.
%
% Every leaf is written on the recto only. Ten views carry text, the even
% ones 322 to 340; the ten odd views 321 to 339 are the blank versos of
% the leaf before, showing only its writing through the paper (view 321
% carries the stamp « Bibliothèque royale »). No view is a second
% exposure; there is no microfilm target.
%
%   v. 322      her § 1: Euler's second solution « sort de l'état de la
%               question », since the motions with the stylet must be among
%               the regular motions of the supported rod, which require
%               sin ω = 0 (« voyez cas IV § 35 »); tang ½ω = (e^ω−1)/(e^ω+1)
%               agrees with it only for ω = 0. « le respect que l'on doit
%               aux opinions d'Euler » makes her treat the general case:
%               § 2 takes up the equation of Euler's § 47 and divides it by
%               e^{2xω} − 1.
%   v. 324      the quotient as a geometric sum when x is an aliquot part;
%               sin ω expanded in powers of sin xω (denominator of x odd or
%               even); sin xω is a factor, whence ω = iπ/x, the rod divided
%               into 1/x equal parts vibrating as separate rods, the
%               « points de repos que l'on pourroit nommer sympatiques »,
%               and the sounds πc√(2gb)/xxaa × 1, 4, 9, 16.
%   v. 326      § 3, Euler's eight coefficients in the general case, γ = −γ';
%               the equations of y for the parts EL, LL', L'L'' … of
%               Euler's fig. 7.
%   v. 328      § 4, the second factor of the final equation cannot vanish
%               except for ω = 0 once sin ω = 0 is imposed; the announcement
%               that for several values of x no solution of Euler's « cas V »
%               (stylet as a fixed point) exists; § 5 begins.
%   v. 330      § 5, the final equation rewritten so as to separate the two
%               portions; the conditions tang xω = (e^{xω}−e^{−xω})/(…) and
%               its mate; the sounds 25, 81, 169, 289 πc√(2gb)/16xxaa of each
%               portion (« Voyez Mémoire d'Euler § 40 »), which can never
%               agree when x' − 1 is even or of the form 4n + 3 (x = 1/x').
%   v. 332      § 6, the conclusion against Euler, and his change of mind
%               between the Methodus inveniendi (« P. 298 § 83 ») and the
%               memoir of 1779 (« § 21 ») on odd numbers of nodes and the
%               « mouvement giratoire »; § 7, the first solution is general.
%   v. 334      end of § 7 (no measurable sound unless the stylet stands at
%               an aliquot division), the analogy with the vibrating string;
%               then « Problème » (§ 8): the rod free at both ends touched by
%               a stylet at any point, « Solution », and Euler's equations
%               I–IV; a boxed marginal nota on square plates, citing
%               Chladni's « traité d'acoustique ».
%   v. 336      equations V–VIII from the free ends (« V. Mémoire d'Euler
%               § 19 »), and the elimination: β', γ' in terms of α'.
%   v. 338      γ', δ' reduced; β, γ, δ in terms of α; substitution into III
%               announced.
%   v. 340      the two members of III computed and equated; equation IV
%               rewritten. The leaf stops there, at the foot of the page:
%               the piece continues past this batch.
%
% The numbers on the leaves, and why no \folio{} is written. Two series,
% one per written leaf, none on the blank versos. (1) At the head, in the
% middle, in parentheses and struck through with a horizontal stroke,
% the page numbers of the piece: (2) at view 322, then (3), (4), (5),
% (6), (7), (8), (9), (10), (11) at views 324 to 340, without a gap ((1)
% is on view 320). (2) At the top right, in the larger, rounder hand and
% broader pen of the running ink series of the earlier batches: 157 at
% view 322, then 158, 159, 160, 161, 162, 163, 164, 165, 166 at views 324
% to 340, one per leaf without a gap (156 is on view 320). Following the
% house rule of this volume, this penned series is recorded here and in
% the notes and is not given as \folio{}; no pencilled foliation was seen.
% Her own paragraph numbers are a third series, written at the head of
% the paragraph: I (v. 322), 2 (v. 322), 3 (v. 326), 4 (v. 328), 5
% (v. 328), 6 (v. 332), 7 (v. 332), 8 (v. 334, the Problème). Other
% numbers are citations: Euler's §§ 19, 21, 35, 40, 47, the Methodus
% p. 298 § 83, « fig. 7 ». Nothing here is Laubenbacher and Pengelley's
% « Manuscript C »: no number theory, and the ink series (157–166) is far
% below 348.
%
% Notation. Hers, kept: x is the fraction of the length at which the
% stylet stands (a letter close to λ on the leaf), a the length, ω the
% angle of the cases, α β γ δ and α' β' γ' δ' Euler's eight coefficients
% (γ written like a « v »), π written like two joined « r », t the time,
% ζ the phase, c, g, b Euler's constants. « sin », « cos », « tang » are
% set \text{sin}, \text{cos}, \text{tang}, with the stop after them where
% it is legible. Braces and brackets are hers. Euler's cases are cited by
% Roman numerals; the I of « IV » is a short stroke joined to the V (views
% 322, 324, 340). Spelling is hers and not flagged: « appuyés »
% (feminine plural), « sympatiques », « suivans », « tems », « fesant »,
% « adaptoit », « aporte », « rejetter », « rejettés », « dela »,
% « delaquestion », « lamême », « demême », « apresent », « àfin »,
% « luimême », « aumilieu », « cequi », « parceque », « parconséquent »,
% « mêmeson », « àlafois », « àdevelopper », « conforment » (for
% « confirment »), « à recours » with the accent.
%
% What could not be read. \ill{} stands once: the last, blotted term of
% the last line at view 336, which the development would make
% cos ω sin (1−x)ω. Readings offered with doubt, \uncertain{}: « fig. » in
% the margin of view 322; « est » for the word the sense wants to be
% « être » at view 322; « =0 » after « sin.ω » at view 324; the exponent
% of L in « L^vi E' » at view 326; « traité » in the marginal nota of
% view 334; « 19 » in « § 19 » at view 336.
%
% At the edges. View 320 (batch 16) opens this piece: title, page (1),
% ink 156, and ends « l'une qui se rapporte au cas […], et l'autre qui
% appartient au cas […] », which view 322 takes up with « la première
% réflexion qui se présente dans l'examen de cette seconde solution ».
% View 340 ends on equation IV rewritten; view 341 is its blank verso; the
% sequel is in batch 18.
%
% Nothing here was checked against any published edition of Euler's memoir
% or of Sophie Germain's papers.
\documentclass[11pt,a4paper]{article}
% The preamble shared by every transcription of the mathematicians' manuscripts in Gallica.
%
% It rests on one idea: **uncertainty compounds, it does not resolve.** A
% transcription that smooths over a doubtful word has destroyed the only thing
% it was made for — a reader must be able to tell, at every point, what was on
% the page and what was supplied. The apparatus macros below are therefore the
% critical apparatus, not decoration.
%
% The same commands are recognised by `scripts/render.mjs`, which produces the
% site's reading view, and by `scripts/tei.mjs`. A macro added here must be
% added there too, or rendering fails loudly — which is the wanted behaviour.
%
% Comments here are English, like the rest of the repository. The documents
% this preamble sets are French, and that is deliberate: see the skills.

\ProvidesPackage{germain}[2026/09/15 Transcriptions of mathematicians' manuscripts in Gallica]

\usepackage{amsmath,amssymb,amsthm}
\usepackage{mathrsfs}
% Kept from the parent preamble so the renderer's diagram support stays
% available; these manuscripts draw no commutative diagrams, and nothing here uses it.
\usepackage{tikz-cd}
\usepackage[english,french]{babel}
\usepackage{fontspec}

% --- Greek ------------------------------------------------------------------
% The text face stays fontspec's default, Latin Modern, which has no Greek:
% XeTeX drops every glyph it lacks with a line in the log and nothing on the
% page, and the Greek words the manuscripts quote vanished from the PDFs. So
% the Greek and Greek Extended blocks get a XeTeX character class of their own,
% and entering or leaving a run of it switches the family to CMU Serif — the
% same Computer Modern design, with polytonic Greek — and back. Only the
% family changes: a Greek word in \textit or \textbf keeps its shape and series.
% The transitions to and from every other class carry the tokens that class
% has towards an ordinary letter, so French spacing before « ; » or inside
% « » still holds after a Greek word. No grouping, so no brace or \endgroup
% can come out unbalanced; maths is left alone, since XeTeX inserts nothing
% there. Kept identical in `exercices.sty`.
\makeatletter
\newfontfamily\germainGreekfont{cmun}[
  Extension=.otf, NFSSFamily=germaingreek,
  UprightFont=*rm, ItalicFont=*ti, SlantedFont=*sl,
  BoldFont=*bx, BoldItalicFont=*bi, BoldSlantedFont=*bl]
\newXeTeXintercharclass\germain@greek
\def\germain@greekrange#1#2{%
  \@tempcnta=#1\relax
  \loop
    \XeTeXcharclass\@tempcnta=\germain@greek
    \ifnum\@tempcnta<#2\advance\@tempcnta\@ne\repeat}
\germain@greekrange{"0370}{"03FF}
\germain@greekrange{"1F00}{"1FFF}
\def\germain@greekprev{\rmdefault}
\def\germain@greekin{%
  \edef\germain@greekprev{\f@family}\fontfamily{germaingreek}\selectfont}
\def\germain@greekout{\fontfamily{\germain@greekprev}\selectfont}
\def\germain@greekjoin{%
  \XeTeXinterchartokenstate=\@ne
  \@tempcnta=\z@
  \loop
    \ifnum\@tempcnta=\germain@greek\else
      \XeTeXinterchartoks\germain@greek\@tempcnta=\expandafter{%
        \expandafter\germain@greekout\the\XeTeXinterchartoks\z@\@tempcnta}%
      \XeTeXinterchartoks\@tempcnta\germain@greek=\expandafter{%
        \the\XeTeXinterchartoks\@tempcnta\z@\germain@greekin}%
    \fi
    \ifnum\@tempcnta<4095\advance\@tempcnta\@ne\repeat}
% babel-french sets its own classes, and saves and restores the interchar
% state, each time a language is selected: join again after it does.
\addto\extrasfrench{\germain@greekjoin}
\addto\extrasenglish{\germain@greekjoin}
\AtBeginDocument{\germain@greekjoin}
\makeatother

\usepackage{geometry}
\geometry{a4paper,margin=25mm,marginparwidth=28mm}
\usepackage{marginnote}
\usepackage{xcolor}
\usepackage[normalem]{ulem}
\setlength{\emergencystretch}{3em}
\usepackage{hyperref}
\hypersetup{colorlinks=true,linkcolor=black,urlcolor=black,pdfborder={0 0 0}}

% --- Batch metadata ---------------------------------------------------------
% Taken as they stand by the rendering script, which shows them at the head of
% the reading view. They fix what the file claims to cover. The macro names are
% the parent project's, so that the scripts stay shared; \folder here means the
% volume's slug — `fr-9115` — and \pages counts Gallica views.

\newcommand{\folder}[1]{\gdef\grothFolder{#1}}
\newcommand{\batch}[1]{\gdef\grothBatch{#1}}
\newcommand{\pages}[2]{\gdef\grothFirst{#1}\gdef\grothLast{#2}}
\newcommand{\foldertitle}[1]{\gdef\grothFoldertitle{#1}}

% The holder's own shelfmark, printed as they write it: « Français 9115 ».
\newcommand{\shelfmark}[1]{\gdef\grothShelfmark{#1}}

% The dating, copied verbatim from the holder's catalogue — never inferred
% here. « XIXe siècle » is the BnF's; a pass that dates a leaf from its content
% says so in a \note{}, as its own inference.
\newcommand{\dating}[1]{\gdef\grothDating{#1}}

% The status notice, printed in the title block and repeated as a diagonal
% watermark on every page of the PDF. The manuscripts are in the public
% domain; what this line declares is not a right but a quality — an
% unverified first pass by a machine. The skills write the line; a file
% without it gets no watermark, so leaving it out is visible in the output.
\newcommand{\watermark}[1]{\gdef\grothWatermark{#1}}

\gdef\grothFolder{?}\gdef\grothBatch{}\gdef\grothFirst{?}\gdef\grothLast{?}%
\gdef\grothFoldertitle{}\gdef\grothShelfmark{}\gdef\grothDating{}\gdef\grothWatermark{}

% --- The résumé -------------------------------------------------------------
\newenvironment{resume}
  {\small\setlength{\parskip}{0.45em}}
  {\par\medskip\hrule\medskip}

% --- Keywords ---------------------------------------------------------------
% In English, closing the modernised reading's résumé: the modern vocabulary
% under which to search for what the leaves construct. The manifest extracts
% them to tag the volume — this line is the single source of the tags.
\newcommand{\keywords}[1]{\par\smallskip\noindent
  {\footnotesize\textsc{Keywords} --- \textit{#1}}\par}

% --- The critical apparatus -------------------------------------------------

% The Gallica view number, in the margin. This is the anchor that lets a
% disagreement over a formula be settled by looking at *one* image rather than
% a volume of 750. The site uses it to turn the facsimile as the transcription
% is scrolled. A view is one scanned image: usually one side of a leaf.
\newcommand{\page}[1]{%
  \marginnote{\footnotesize\color{gray!70}\textsf{v.\,#1}}%
  \label{page:#1}\ignorespaces}

% The BnF's pencilled foliation, where the leaf carries it and a pass can read
% it: « 198r ». This is what the literature cites — « ff. 198r–208v » — and
% the one line that turns a citation into a view. Recorded, never inferred:
% a folio a pass cannot read is not written.
\newcommand{\folio}[1]{%
  \marginnote{\footnotesize\color{gray!50}\textsf{f.\,#1}}\ignorespaces}

% The modernised reading's coarser counterpart to \page{}: which views a
% section is drawn from.
\newcommand{\pagerange}[2]{%
  \marginnote{\footnotesize\color{gray!70}\textsf{v.\,#1--#2}}%
  \ignorespaces}

% Illegible: never guessed.
\newcommand{\ill}{\textcolor{red!60!black}{[\,\ldots\,]}}

% A doubtful reading: the word is given, and flagged as uncertain.
\newcommand{\uncertain}[1]{\uline{#1}}

% An editorial addition — a missing letter, an implied word.
\newcommand{\add}[1]{\textcolor{blue!50!black}{[#1]}}

% Struck out by the writer. What was crossed out is part of the document.
\newcommand{\struck}[1]{\sout{#1}}

% The transcriber's note, on the state of the page or the mathematics.
\newcommand{\note}[1]{\par\smallskip{\small\color{gray!60!black}\textit{[#1]}}\par\smallskip}

% A marginal note of hers, distinct from the body of the text.
\newcommand{\marginal}[1]{\par\smallskip\noindent\hspace{1em}%
  {\small\color{gray!60!black}#1}\par\smallskip}

% --- Title ------------------------------------------------------------------
% The title block is French, like the documents themselves.

\AddToHook{shipout/background}{%
  \ifx\grothWatermark\empty\else
    \begin{tikzpicture}[remember picture, overlay]
      \node[rotate=52, scale=3.4, align=center, text=gray!18,
            font=\sffamily\bfseries]
        at (current page.center) {\grothWatermark};
    \end{tikzpicture}%
  \fi}

\AtBeginDocument{%
  \begin{center}
    {\large\bfseries Manuscrits de mathématiciens (Gallica) ---
      \ifx\grothShelfmark\empty\grothFolder\else\grothShelfmark\fi}\\[2pt]
    {\itshape\grothFoldertitle}\\[4pt]
    {\small vues \grothFirst--\grothLast
      \ifx\grothBatch\empty\else\ (lot \grothBatch)\fi}\\[2pt]
    \ifx\grothDating\empty\else
      {\small\color{gray!60!black}Datation du catalogue : \grothDating}\\[2pt]
    \fi
    \ifx\grothWatermark\empty\else
      {\small\bfseries\color{red!45!gray}\grothWatermark}\\[2pt]
    \fi
    {\footnotesize\color{gray!60!black}Transcription établie d'après les images IIIF de Gallica.
     Source gallica.bnf.fr / Bibliothèque nationale de France.\\
     Les lectures incertaines sont soulignées, les illisibles notées [\,\ldots\,],
     les ajouts entre crochets ; \textsf{v.} numérote les vues de Gallica,
     \textsf{f.} la foliotation au crayon de la BnF.}
  \end{center}
  \vspace{1em}}

\endinput


\folder{fr-9115}
\shelfmark{Français 9115}
\batch{17}
\pages{321}{340}
\dating{1801-1900}
\watermark{Lecture automatique\\non vérifiée}
\foldertitle{Papiers de Sophie GERMAIN. — Recueil de dissertations et problèmes mathématiques et physiques.}

\begin{document}

\note{Numérotation des feuillets. Chaque feuillet écrit porte deux nombres,
et les versos blancs n'en portent aucun : au milieu de la tête, entre
parenthèses et barré d'un trait, le numéro de page de la pièce, (2) à la
vue 322, puis (3) à (11) aux vues 324 à 340 ; en haut à droite, d'une
plume plus large et d'une écriture plus ronde, 157 à la vue 322, puis 158
à 166 aux vues 324 à 340, un par feuillet et sans lacune. Cette seconde
série, tracée à la plume, prolonge la série à l'encre relevée dans les
lots précédents de ce volume ; aucune foliotation au crayon n'a été lue,
et aucun « f. » n'est donc porté. Tous ces nombres ont été lus sur la
vue ; aucun n'est déduit. Les vues 321, 323, 325, 327, 329, 331, 333,
335, 337 et 339 sont les versos blancs des feuillets qui les précèdent et
ne sont pas transcrites ; la vue 321 porte le timbre « Bibliothèque
royale ». La pièce commence à la vue 320 (lot 16), sous le titre
« Remarques sur le mémoire d'Euler », et se poursuit au-delà de la vue
340.}

\page{322}

\marginal{\uncertain{fig.} 7}
I La première réflexion qui se présente dans l'examen de cette seconde solution, est
qu'elle sort de l'état delaquestion; car l'énoncé du problème indique assez que les
mouvemens dela lame soumise à l'action du stilet en L, doivent être compris au
nombre detous les mouvemens réguliers possibles dont la même lame est susceptible
lorsqu'elle a ses extremités appuyés; et les formules employés à trouver l'équation
finale, conforment cette remarque puisque l'auteur à recours aquatre équations
dépendantes de l'état des extremités. Cela posé, il en résulte, ce me semble,
que la condition $\text{sin.}\,\omega = 0$ qu'il a établi devoir toujours être remplie pour tous les
mouvemens réguliers possibles dela lame dont il s'agit (voyez cas IV §35), doit
s'accorder avec la solution du présent problème; et c'est effectivement cequ'on
trouve pour la première solution; savoir, $\text{sin}\,\frac{1}{2}\omega = 0$: mais lamême chose
n'a pas lieu pour la seconde; car l'équation $\text{tang}\,\frac{1}{2}\omega = \frac{e^{\omega}-1}{e^{\omega}+1}$ ne peut
s'accorder avec la condition $\text{sin.}\,\omega = 0$, que dans la seule supposition $\omega = 0$.
Cependant le respect que l'on doit aux opinions d'Euler, me détermine a
examiner le cas général et a chercher s'il est susceptible d'une solution
dans laquelle celle qui m'a inspiré des doutes, puisse \uncertain{est} comprise.

2. Pour cela, je reprens l'équation (§ 47 d'Euler; savoir
\[ 0 = 2(1-e^{2\omega}) - (e^{x\omega} - e^{(2-x)\omega})(e^{x\omega} - e^{-x\omega})\,\frac{\text{sin}\,\omega}{\text{sin}\,x\omega\,\text{sin}\,(1-x)\omega} \]
ou
\[ 0 = 2(1-e^{2\omega})\,\text{sin}\,x\omega\,\text{sin}\,(1-x)\omega - (1-e^{2(1-x)\omega})(e^{2x\omega}-1)\,\text{sin}\,\omega \]
Si on veut conserver l'analogie entre le cas général et celui ou
$x = \frac{1}{2}$, on peut diviser cette équation par le facteur $e^{2x\omega}-1$, duquel on
tire, comme dans le cas cité, $\omega = 0$; et elle se réduit par cette division à
\note{En tête, au milieu, « (2) » biffé d'un trait ; en haut à droite
« 157 », à l'encre. Main de mise au net, régulière, presque sans
rature. Dans la marge de gauche, en face de la troisième ligne, « fig. 7 »
(lecture de la première syllabe incertaine). Le chiffre I en tête du
premier alinéa et le 2 du second sont ses numéros de paragraphe.
« cas IV §35 » : le I n'est qu'un trait court collé au V, et le signe
de paragraphe, tracé petit, se lit aussi « s ». Le § 35 d'Euler est le
problème de son cas IV, les deux extrémités appuyées, copié en latin aux
vues 181 à 183 de ce volume. À la dernière ligne du premier paragraphe, le
mot entre « puisse » et « comprise » se lit « est » ; le sens demande
« être ». La parenthèse ouverte avant « § 47 » n'est pas fermée. La
variable $x$, fraction de la longueur, est tracée d'une lettre proche
d'un $\lambda$. Le verso (vue 323) est blanc.}

\page{324}

\[ 2\{1 + e^{2x\omega} + e^{4x\omega} + e^{6x\omega} \ldots\ldots + e^{2(1-x)\omega}\}\,\text{sin}\,x\omega\,\text{sin}\,(1-x)\omega + (1 - e^{2(1-x)\omega})\,\text{sin.}\,\omega = 0 \]
J'observe que, par la nature delaquestion, $x$ est toujours égal à une fraction, et
que, dans le cas des deux extrémités appuyés, cas dont il s'agit ici, cette fraction
est partie aliquote de l'unité, ou multiple d'une telle partie.

Lorsque le numérateur est l'unité, on trouve par les formules connues, d'abord,
pour le cas où le dénominateur dela même fraction $x$ est un nombre impair,
\[ \text{sin.}\,\omega = \frac{1}{x}\,\text{sin.}\,x\omega - \frac{1-x^2}{x^3.2.3}\,\text{sin}^3 x\omega + \frac{(1-x^2)(1-9x^2)}{x^5.2.3.4.5}\,\text{sin}^5 x\omega\ \&c., \]
et ensuite, pour le cas ou le dénominateur de cette fraction $x$ est un nombre pair,
\[ \text{sin.}\,\omega = \text{cos}\,x\omega\left\{\frac{1}{x}\,\text{sin}\,x\omega - \frac{1-4x^2}{x^3\,2.3}\,\text{sin}^3 x\omega + \frac{(1-4x^2)(1-16x^2)}{x^5.2.3.4.5}\,\text{sin}^5 x\omega - \&c\right\}. \]
L'équation qu'il s'agit de résoudre, se réduit donc à l'une ou l'autre des formes
suivantes
\[ 2\{1 + e^{2x\omega} + e^{4x\omega} + e^{6x\omega} \ldots + e^{2(1-x)\omega}\}\,\text{sin}\,x\omega\,\text{sin}\,(1-x)\omega \]
\[ + (1-e^{2(1-x)\omega})\left\{\frac{1}{x} - \frac{1-x^2}{x^3.2.3}\,\text{sin}^2 x\omega + \frac{(1-x^2)(1-9x^2)}{x^5.2.3.4.5}\,\text{sin.}^4 x\omega - \&c\right\}\text{sin}\,x\omega = 0 \]
\[ 2\{1 + e^{2x\omega} + e^{4x\omega} + e^{6x\omega} \ldots + e^{2(1-x)\omega}\}\,\text{sin.}\,x\omega\,\text{sin.}\,(1-x)\omega \]
\[ + (1-e^{2(1-x)\omega})\,\text{cos}\,x\omega\left\{\frac{1}{x} - \frac{1-4x^2}{x^3.2.3}\,\text{sin}^2 x\omega + \frac{(1-4x^2)(1-16x^2)}{x^5.2.3.4.5}\,\text{sin}^4 x\omega\right\}\text{sin}\,x\omega = 0 \]
ainsi, dans l'un et l'autre cas, $\text{sin}\,x\omega$ est un des facteurs de cette équation; et, en
égalant ce facteur à zéro, on a une solution qui se rapporte au cas IV et de
laquelle il résulte, comme dans la supposition $x = \frac{1}{2}$, $\omega = \frac{i\pi}{x}$, $i$ étant un nombre
entier quelconque. on en conclut encore, comme dans cette supposition, que la
lame se partage alors en $\frac{1}{x}$ parties de longueur $xa$, et que ces différentes
parties font leurs vibrations de la même manière que si elles étoient autant de
lames séparées et de pareille longueur, qui auroient leurs extrémités appuyés, tant
dans le point où est placé le stilet, que dans les autres points de repos que
l'on pourroit nommer sympatiques. Enfin il en résulte aussi que les sons
simples que les parties fractionnaires dela lame peuvent donner, sont exprimés
par les nombres suivans: $\frac{\pi c\sqrt{2gb}}{xxaa}$; $\frac{4\pi c\sqrt{2gb}}{xxaa}$; $9\frac{\pi c\sqrt{2gb}}{xxaa}$; $\frac{16\pi c\sqrt{2gb}}{xxaa}$; \&c.

Il est clair que la condition $\text{sin.}\,\omega$ \uncertain{$=0$}, qui, par la nature dela question, doit avoir lieu
en même tems que celle qui résulte dela solution de l'équation du problème,
\note{En tête, au milieu, « (3) » biffé ; en haut à droite « 158 », à
l'encre. La page s'ouvre sur l'équation qui achève la phrase de la vue
322 (« elle se réduit par cette division à »). Dans la première des deux
« formes suivantes », le premier facteur est écrit sans « cos », dans la
seconde avec « $\text{cos}\,x\omega$ » : c'est la différence des deux
développements de $\text{sin.}\,\omega$ qui précèdent ; la seconde forme
n'a pas de « \&c » dans son accolade. Le $\pi$ est tracé comme deux « r »
liés. Après « la condition $\text{sin.}\,\omega$ », à l'avant-dernière
ligne, un petit signe en indice, qui ressemble à une croix, tient
apparemment lieu de « $=0$ » ; lecture incertaine. « cas IV » : le I est un trait court collé au V, comme à la vue
322. « dénominateur » est
écrit deux fois avec un « n » repassé. Le verso (vue 325) est blanc.}

\page{326}

est satisfaite ici, puisqu'elle résulte dela supposition $\text{sin}\,x\omega = 0$.

3. Cela posé, on peut encore suivre, dans la détermination des coëfficiens
$\alpha, \beta, \gamma, \delta$ et $\alpha', \beta', \gamma', \delta'$, l'analogie entre le cas général et celui qu'Euler a calculé;
car on trouve delamême maniere
\[ \alpha = e^{x\omega} - e^{(2-x)\omega};\quad \beta = e^{(2-x)\omega} - e^{x\omega};\quad \gamma = \frac{(e^{(2-x)\omega} - e^{x\omega})(e^{x\omega} - e^{-x\omega})}{\text{sin}\,x\omega};\quad \delta = 0 \]
\[ \alpha' = e^{x\omega} - e^{-x\omega};\quad \beta' = -e^{2\omega}(e^{x\omega} - e^{-x\omega}); \]
\[ \gamma' = \frac{(e^{(2-x)\omega} - e^{x\omega})(e^{x\omega} - e^{-x\omega})}{\text{sin.}\,x\omega - \text{cos}\,x\omega\,\text{tang.}\,\omega}; \]
\[ \delta' = \frac{(e^{x\omega} - e^{(2-x)\omega})(e^{x\omega} - e^{-x\omega})\,\text{tang}\,\omega}{\text{sin.}\,x\omega - \text{cos}\,x\omega\,\text{tang.}\,\omega}; \]
et, à cause de $\text{cos.}\,\omega = 1$ d'ou il résulte $\text{sin.}\,x\omega - \text{cos}\,x\omega\,\text{tang}\,\omega = -\,\text{sin}\,(1-x)\omega$, on peut écrire
\[ \gamma' = -\frac{(e^{(2-x)\omega} - e^{x\omega})(e^{x\omega} - e^{-x\omega})}{\text{sin}\,(1-x)\omega}\qquad \delta' = \frac{(e^{x\omega} - e^{(2-x)\omega})(e^{x\omega} - e^{-x\omega})}{\text{sin}\,(1-x)\omega}. \]
Si on multiplie donc les premiers coëfficiens par $\text{sin}\,x\omega$, et les seconds par $\text{sin}\,(1-x)\omega$,
cequi revient au même puisque ces deux sinus sont l'un et l'autre égaux
a zéro, on aura, comme dans le cas ou $x = \frac{1}{2}$,
\[ \alpha = 0,\quad \beta = 0,\quad \gamma = -(e^{x\omega} - e^{(2-x)\omega})(e^{x\omega} - e^{-x\omega}),\quad \delta = 0 \]
\[ \alpha' = 0,\quad \beta' = 0,\quad \gamma' = (e^{x\omega} - e^{(2-x)\omega})(e^{x\omega} - e^{-x\omega}),\quad \delta' = 0 \]
parconséquent $\gamma = -\gamma'$; et on peut faire également ici $\gamma = 1$ $\gamma' = -1$

Ainsi, en prenant pour exemple le cas représenté par la fig. 7, et
supposant le stilet placé en L, l'équation qui déterminera le mouvement
de la partie EL, sera $y = c\,\text{sin}\,(\zeta + t.\frac{\omega\omega c\sqrt{2gb}}{aa})\,\text{sin}\,x\omega$; et, pour la partie
L\uncertain{$^{\text{vi}}$}E$'$, en tout semblable à la partie LL$'$, on aura $y = -c\,\text{sin}\,(\zeta + t.\frac{\omega\omega c\sqrt{2gb}}{aa})\,\text{sin}\,x\omega$;
il est clair que le mouvement dela partie L$'$L$''$ est déterminé demême
par l'équation $y = c\,\text{sin}\,(\zeta + t.\frac{\omega\omega c\sqrt{2gb}}{aa})\,\text{sin}\,x\omega$, celui de la partie L$''$L$'''$ par
l'équation $y = -c\,\text{sin}\,(\zeta + t.\frac{\omega\omega c\sqrt{2gb}}{aa})\,\text{sin}\,x\omega$; et ainsi de suite, en adoptant
le signe $+$ pour les parties dela courbe qui sont situés audessus
de l'axe, et le signe $-$ pour celles qui sont placés au dessous
du même axe, ou ligne de repos.
\note{En tête, au milieu, « (4) » biffé ; en haut à droite « 159 », à
l'encre. Le premier mot, « est satisfaite », achève la phrase laissée
ouverte au bas de la vue 324. Le mot « cas », dans « le cas général »,
est repassé et empâté, non biffé ; la dernière lettre de « calculé » est
de même surchargée. Au signe d'égalité de la première valeur de
$\gamma'$ s'ajoute un petit point en dessous, à droite, qui n'a pas été lu
comme un signe. Dans la seconde série de valeurs, après « $\gamma =$ »,
un trait noir épais et empâté tient la place d'un signe moins, comme le
demande « $\gamma = -\gamma'$ » à la ligne suivante ; il est donné comme
tel. L'exposant de L dans « L$^{\text{vi}}$E$'$ » est d'une lecture
incertaine. Dans les quatre équations de $y$, le dernier facteur est
écrit « sin » suivi de la même lettre que la fraction $x$ et de
$\omega$ ; il est transcrit ainsi. Le verso (vue 327) est blanc.}

\page{328}

4. Je reviens à présent à l'équation du problème: elle a pour second facteur
\[ 2\{1 + e^{2x\omega} + e^{4x\omega} + e^{6x\omega} \ldots + e^{2(1-x)\omega}\}\,\text{sin}\,(1-x)\omega \]
\[ + (1-e^{2(1-x)\omega})\left\{\frac{1}{x} - \frac{1-x^2}{2.3.x^3}\,\text{sin.}^2 x\omega + \frac{(1-x^2)(1-9x^2)}{2.3.4.5.x^5}\,\text{sin}^4 x\omega - \&c\right\} \]
ou
\[ 2\{1 + e^{2x\omega} + e^{4x\omega} + e^{6x\omega} \ldots + e^{2(1-x)\omega}\}\,\text{sin}\,(1-x)\omega \]
\[ + (1-e^{2(1-x)\omega})\left\{\frac{1}{x} - \frac{1-4x^2}{2.3.x^3}\,\text{sin.}^2 x\omega + \frac{(1-4x^2)(1-16x^2)}{2.3.4.5.x^5}\,\text{sin}^4 x\omega - \&c\right\}\text{cos}\,x\omega \]
suivant que le dénominateur de la fraction $x$ est impair ou pair.

D'après les valeurs de $\text{sin.}\,\omega$ rapportés cidessus N.o 2, et en admettant que la
condition $\text{sin.}\,\omega = 0$ doive s'accorder avec toutes les solutions de ce problème, on voit
qu'il faut nécessairement que $\text{sin.}\,x\omega = 0$, ou que ce soit le second facteur de $\text{sin.}\,\omega$
savoir: $\frac{1}{x} - \frac{1-x^2}{2.3.x^3}\,\text{sin.}^2 x\omega + \&c.$ ou $\left\{\frac{1}{x} - \frac{1-4x^2}{2.3.x^3}\,\text{sin.}^2 x\omega + \&c.\right\}\text{cos}\,x\omega$ suivant les cas, qui soit
egal a zéro. mais, si on adaptoit cette dernière supposition, le second facteur
de l'équation du problème se réduiroit, dans tous les cas, à $2\{1 + e^{2x\omega} + e^{4x\omega} \ldots + e^{2(1-x)\omega}\}\,\text{sin}\,(1-x)\omega$
quantité qui ne peut être egalé a zéro, qu'en fesant $\text{sin}\,(1-x)\omega = 0$. ainsi, a cause
de $\text{sin.}\,\omega = 0$, on seroit ramené àla condition $\text{sin.}\,x\omega = 0$. En adoptant donc cette
dernière hypothèse, le second facteur dont il s'agit, se trouve réduit à $(1-e^{2(1-x)\omega})\frac{1}{x}$;
et il est visible que cette quantité ne peut être egalé à zéro, que dans la
seule supposition $\omega = 0$. J'ai déjà fait cette remarque pour le cas ou $x = \frac{1}{2}$: mais
je vais prouver, deplus, que, dans le cas général, non seulement \struck{non seulement}
le second facteur de l'équation du problème ne peut être egalé a zéro que
dans l'hypothèse $\omega = 0$ qui se rapporte au cas du repos dela lame, lorsqu'on a
égard a la condition $\text{sin}\,\omega = 0$, mais encore qu'indépendamment de cette condition,
il se trouve plusieurs valeurs de $x$, pour lesquelles l'équation qu'il s'agit de
résoudre, n'est susceptible d'aucune solution qui puisse se rapporter au cas V,
c'est-à-dire qui puisse mener a considérer les deux portions dela lame
séparés par le stilet comme fixés a ce point et appuyés a leur autre
extrémité, comme Euler le suppose dans sa seconde solution.

5 Pour cela, je reprens l'équation finale
\note{En tête, au milieu, « (5) » biffé ; en haut à droite « 160 », à
l'encre. Un trait vertical fin traverse la première grande formule sans
en toucher la lecture. Dans « N.o 2 », le N est surchargé et empâté ; le
renvoi est à son paragraphe 2 (vues 322--324), où sont données les
valeurs de $\text{sin.}\,\omega$. « adaptoit » est écrit ainsi. Le « non
seulement » répété à la fin d'une ligne est biffé. Le chiffre après « au
cas » est un V dont le premier jambage est empâté, qui ressemble à un
4 ; il répond au « cas V » nettement écrit de la vue 330, où revient la
même idée (le stilet regardé comme un point fixe, l'autre extrémité
appuyée).
Le « r » d'« Euler » est surchargé. La page s'arrête sur « l'équation
finale », au tiers inférieur du feuillet, le reste restant blanc ; la
suite est à la vue 330. Le verso (vue 329) est blanc.}

\page{330}

\[ 0 = 2(1-e^{2\omega}) - (e^{x\omega} - e^{(2-x)\omega})(e^{x\omega} - e^{-x\omega})\,\frac{\text{sin.}\,\omega}{\text{sin}\,x\omega\,\text{sin}\,(1-x)\omega} \]
et je la transforme comme il suit
\[ 0 = \frac{2(e^{\omega} - e^{-\omega})}{(e^{(x-1)\omega} - e^{(1-x)\omega})(e^{x\omega} - e^{-x\omega})}\,.\,\text{sin}\,x\omega\,\text{sin}\,(1-x)\omega - \text{sin}\,\omega \]
\[ 0 = \frac{\begin{array}{c}(e^{(x-1)\omega} + e^{(1-x)\omega})(e^{x\omega} - e^{-x\omega}) \\  + (e^{x\omega} + e^{-x\omega})(e^{(1-x)\omega} - e^{(x-1)\omega})\end{array}}{(e^{(1-x)\omega} - e^{(x-1)\omega})(e^{x\omega} - e^{-x\omega})}\,\text{sin.}\,x\omega\,\text{sin.}\,(1-x)\omega \]
\[ - \{\text{sin}\,x\omega\,\text{cos}\,(1-x)\omega + \text{sin}\,(1-x)\omega\,\text{cos}\,x\omega\} \]
\[ 0 = \frac{e^{(x-1)\omega} + e^{(1-x)\omega}}{e^{(1-x)\omega} - e^{(x-1)\omega}}\,\text{sin.}\,x\omega\,\text{sin.}\,(1-x)\omega - \text{cos}\,(1-x)\omega\,\text{sin}\,x\omega \]
\[ + \frac{e^{x\omega} + e^{-x\omega}}{e^{x\omega} - e^{-x\omega}}\,\text{sin}\,x\omega\,\text{sin}\,(1-x)\omega - \text{sin}\,(1-x)\omega\,\text{cos}\,x\omega \]
et enfin
\[ 0 = \text{sin}\,x\omega\left\{\frac{e^{(1-x)\omega} + e^{(x-1)\omega}}{e^{(1-x)\omega} - e^{(x-1)\omega}}\,\text{sin.}\,(1-x)\omega - \text{cos.}\,(1-x)\omega\right\} \]
\[ + \text{sin}\,(1-x)\omega\left\{\frac{e^{x\omega} + e^{-x\omega}}{e^{x\omega} - e^{-x\omega}}\,\text{sin.}\,x\omega - \text{cos.}\,x\omega\right\} \]
si on pouvoit considérer le stilet placé en L comme fixé, il faudroit que
les deux conditions $\frac{e^{(1-x)\omega} + e^{(x-1)\omega}}{e^{(1-x)\omega} - e^{(x-1)\omega}}\,\text{sin}\,(1-x)\omega = \text{cos}\,(1-x)\omega$ et $\frac{e^{x\omega} + e^{-x\omega}}{e^{x\omega} - e^{-x\omega}}\,\text{sin}\,x\omega = \text{cos.}\,x\omega$
ou $\text{tang.}\,(1-x)\omega = \frac{e^{(1-x)\omega} - e^{(x-1)\omega}}{e^{(1-x)\omega} + e^{(x-1)\omega}}$ et $\text{tang.}\,x\omega = \frac{e^{x\omega} - e^{-x\omega}}{e^{x\omega} + e^{-x\omega}}$ eussent lieu
en même tems; et on auroit alors une solution qui pourroit être rapportée
au cas V: mais on verra aisément que, dans la plupart des cas, ces deux
conditions ne peuvent exister àlafois pour aucune autre supposition que celle $\omega = 0$.
En effet, il faut au moins que les deux portions dela lame rendent le
mêmeson; car cette supposition est fondamentale. Or, si l'équation
$\text{tang}\,x\omega = \frac{e^{x\omega} - e^{-x\omega}}{e^{x\omega} + e^{-x\omega}}$ a lieu pour la portion EL de longueur $xa$, les sons
simples que cette portion pourra rendre seront exprimés par les nombres
suivans: $\frac{25\pi c\sqrt{2gb}}{16xxaa}$; $\frac{81\pi c\sqrt{2gb}}{16xxaa}$; $\frac{169\pi c\sqrt{2gb}}{16xxaa}$; $\frac{289\pi c\sqrt{2gb}}{16xxaa}$; \&c. (Voyez Memoire d'Euler § 40);
et, en supposant en même tems que l'autre équation $\text{tang}\,(1-x)\omega = \frac{e^{(1-x)\omega} - e^{(x-1)\omega}}{e^{(1-x)\omega} + e^{(x-1)\omega}}$
ait pareillement lieu pour la portion LF de longueur $(1-x)a$, les sons simples
que cette portion pourra rendre, seront exprimés par les nombres
\[ \frac{25\pi c\sqrt{2gb}}{16(1-x)^2 aa};\quad \frac{81\pi c\sqrt{2gb}}{16(1-x)^2 aa};\quad \frac{169\pi c\sqrt{2gb}}{16(1-x)^2 aa};\quad \frac{289\pi c\sqrt{2gb}}{16(1-x)^2 aa},\ \&c. \]
et il est clair, en fesant pour plus desimplicité $x = \frac{1}{x'}$ et parconséquent $(1-x)^2 = \frac{(1-x')^2}{x'x'}$,
qu'aucun nombre de la première suite ne pourra être égal aux nombre quelconques
dela seconde, lorsque $x'-1$, sera un nombre pair ou un nombre impair delaforme $4n+3$.
\note{En tête, au milieu, « (6) » biffé ; en haut à droite « 161 », à
l'encre. La page s'ouvre sur « l'équation finale » annoncée au bas de la
vue 328, qui est celle du § 47 d'Euler déjà écrite à la vue 322. Dans
les deux dénominateurs de la troisième et de la quatrième équation, les
exposants sont surchargés : « $(1-x)\omega$ » et « $(x-1)\omega$ » sont
écrits sur l'ordre inverse ; la leçon finale est donnée. Dans la
deuxième équation, le dénominateur garde l'ordre
$e^{(x-1)\omega} - e^{(1-x)\omega}$. L'extrémité de la lame opposée à E
est ici appelée F (« la portion LF »), la lettre étant nettement formée.
Le verso (vue 331) est blanc.}

\page{332}

6. J'ai voulu entrer dans tous ces détails, àfin d'être autorisée a conclure qu'Euler
a été induit en erreur, lorsqu'il a cru pouvoir admettre sa seconde solution
dans laquelle il suppose seulement que les deux portions EL, LE$'$ puissent se mouvoir
en même tems; et c'est pourquoi j'ai fait voir que cette condition même ne peut
pas être remplie pour toutes les valeurs possibles de $x$. D'ailleurs, quelque
répugnance que l'on puisse avoir à devier de l'opinion d'Euler, on est pourtant
forcé de convenir que, sans sortir du sujet de ce mémoire, il a luimême varié
de sentiment. En effet, quoique, dans sa \emph{methodus inveniendi curvas \&c},
il ait suivi, pour l'analyse du cas où les deux extrémités de la lame élastique
vibrante sont libres, une méthode qui ne diffère pas essentiellement de celle qu'il
a depuis adopté dans son mémoire de 1779, il s'est pourtant décidé, dans son
premier travail, à rejetter toutes les solutions qui donnent lieu à un nombre
impair de nœuds. la raison qu'il en aporte, est qu'il existeroit alors un
mouvement giratoire dela lame autour du nœud placé aumilieu de cette même lame
(V. methodus \&c. P. 298 § 83); et, au contraire, il a admis dans son mémoire les solutions qu'il
avoit d'abord rejettés, et a simplement employé la considération du mouvement
giratoire, à prouver que, dans le mouvement le plus simple de la lame élastique
vibrante dont les deux extrémités sont libres, la courbe doit couper l'axe en deux
points (V. Memoire § 21).

7. Je reviens encore à la première solution d'Euler. on voit qu'elle est générale,
et qu'elle montre comment il s'établit un nombre de nœuds dépendant dela
valeur de $x$; car, puisqu'on a $y = 0$ par la supposition $\text{sin}\,x\omega = 0$, on aura
encore $y = 0$, pour les points qui correspondent aux sinus multiples tels que
$\text{sin.}\,2x\omega$ $\text{sin.}\,3x\omega$, \&c.: d'où il résulte que le mouvement dela lame est lemême,
soit qu'on place \add{le} stilet \struck{soit placé} a l'un ou l'autre de ces points, pourvu toutefois
qu'en désignant par $\text{sin}\,nx\omega$ un des sinus de cette suite, $n$ ne soit pas facteur du
dénominateur de la fraction $x$; car alors la supposition $\text{sin}\,nx\omega = 0$, n'entraîneroit pas nécessairement $\text{sin}\,x\omega = 0$.
\note{En tête, au milieu, « (7) » biffé ; en haut à droite « 162 », à
l'encre. « \emph{methodus inveniendi curvas \&c} » est souligné sur le
feuillet. Plusieurs mots sont surchargés sans être biffés : « tems »,
« pas » en tête de ligne, « D'ailleurs », « l'opinion ». Dans « soit
qu'on place le stilet », « qu'on » est écrit sur une autre leçon et
« le » est ajouté au-dessus de la ligne ; « soit placé », qui suivait
« stilet », est biffé. La dernière ligne est serrée contre le bord et
son dernier « $=0$ » est à peine formé. Le reste du feuillet est blanc.
Le verso (vue 333) est blanc.}

\page{334}

Il résulte encore de cette solution que tout mouvement régulier, et parconséquent
tout son mesurable, est impossible à obtenir de la lame dont les deux extrémités
sont appuyés, lorsque le stilet est placé à tout autre point qu'à une division
aliquote de la longueur de la même lame. on voit en effet qu'excepté
ce cas, les conditions $\text{sin.}\,\omega = 0$ et $\text{sin.}\,x\omega = 0$ n'ont pas lieu en même tems,
et que l'équation finale du problème n'a plus $\text{sin}\,x\omega$ pour facteur.

Ces observations établissent les plus grands rapports entre les circonstances
du mouvement d'une telle lame et celles qui ont lieu pour la corde vibrante:
mais malheureusement ce cas des deux extrémités appuyés est le seul où de
semblables analogies se fassent remarquer

Je me propose apresent d'appliquer la methode d'Euler au cas où le stilet
est placé dans un point quelconque de la longueur d'une lame élastique
vibrante dont les deux extrémités sont libres. L'importance de ce cas$^{\times}$ m'engage
àdevelopper ici les calculs qui y sont relatifs.
\marginal{$\times$ N.ta Les vibrations transversales des plaques élastiques carrées y sont comprises, puisque la largeur n'influe pas (V. \uncertain{traité} d'acoustique de Chladni)}

\section{Problème}

8. Si une lame élastique, libre à l'une et à l'autre de ses extrémités,
est soumise à l'action d'un stilet dans un point quelconque de sa longueur:
trouver tous les mouvemens de vibration dont elle est susceptible?

\subsection{Solution}

Je reprens les équations
\begin{align*}
&\text{I} && \alpha e^{x\omega} + \beta e^{-x\omega} + \gamma\,\text{sin.}\,x\omega + \delta\,\text{cos.}\,x\omega = 0\\
&\text{II.} && \alpha' e^{x\omega} + \beta' e^{-x\omega} + \gamma'\,\text{sin.}\,x\omega + \delta'\,\text{cos}\,x\omega = 0\\
&\text{III} && \alpha e^{x\omega} - \beta e^{-x\omega} + \gamma\,\text{cos}\,x\omega - \delta\,\text{sin}\,x\omega = \alpha' e^{x\omega} - \beta' e^{-x\omega} + \gamma'\,\text{cos}\,x\omega - \delta'\,\text{sin.}\,x\omega\\
&\text{IV} && (\alpha-\alpha')e^{x\omega} + (\beta-\beta')e^{-x\omega} - (\gamma-\gamma')\,\text{sin}\,x\omega - (\delta-\delta')\,\text{cos}\,x\omega = 0
\end{align*}
données par Euler dans l'analyse du problème semblable relatif au cas où les
extrémités dela lame sont l'une et l'autre appuyés \add{(voyez §47)}, parcequ'il est clair que ces équations
sont également applicables au cas présent, puisqu'elles sont indépendantes de l'état des
extrémités.
\note{En tête, au milieu, « (8) » biffé ; en haut à droite « 163 », à
l'encre. Le premier alinéa achève le paragraphe 7 de la vue 332. La note
marginale, appelée par une croix après « ce cas », est écrite dans la
marge de gauche, encadrée d'un trait, et commence par « N.ta » (nota) ;
le mot « traité » y est surchargé. « Problème » et « Solution » sont
centrés et soulignés sur le feuillet. « (voyez §47) » est ajouté
au-dessus de la ligne, après « appuyés » ; le § 47 d'Euler, la lame
fixée en un point intermédiaire L, est copié en latin à partir de la vue
192. Le verso (vue 335) est blanc.}

\page{336}

Maintenant donc, en fesant successivement $x = 0$ et $x = 1$, il résultera quatre équations
de la considération des extrémités savoir: (V. Memoire d'Euler § \uncertain{19})
\[ \text{V.}\ \alpha - \beta = \gamma,\qquad \text{VI}\ \alpha + \beta = \delta \]
\[ \text{VII}\ \alpha' e^{\omega} + \beta' e^{-\omega} - \gamma'\,\text{sin.}\,\omega - \delta'\,\text{cos}\,\omega = 0 \quad\text{et}\quad \text{VIII}\ \alpha' e^{\omega} - \beta' e^{-\omega} - \gamma'\,\text{cos.}\,\omega + \delta'\,\text{sin.}\,\omega = 0 \]
ainsi on aura, comme dans le cas des extrémités appuyés, huit équations pour
déterminer les huit coëfficiens $\alpha, \beta, \gamma, \delta$ et $\alpha', \beta', \gamma', \delta'$; et l'on parviendra demême
à une équation finale, de laquelle on pourra tirer la valeur d'$\omega$.

Pour cela, je commence par multiplier l'équation VII par $\text{sin.}\,\omega$, et l'équation VIII
par $\text{cos.}\,\omega$; puis je les ajoute l'une a l'autre, cequi donne
\[ \alpha' e^{\omega}(\text{sin.}\,\omega + \text{cos.}\,\omega) + \beta' e^{-\omega}(\text{sin}\,\omega - \text{cos.}\,\omega) = \gamma' \]
je multiplie ensuite l'équation VII par $\text{cos.}\,\omega$, et l'équation VIII par $\text{sin.}\,\omega$; cequi donne,
en les retranchant l'une de l'autre,
\[ \alpha' e^{\omega}(\text{cos.}\,\omega - \text{sin.}\,\omega) + \beta' e^{-\omega}(\text{sin.}\,\omega + \text{cos.}\,\omega) = \delta' \]
En substituant les présentes valeurs de $\gamma'$ et $\delta'$ dans l'équation II, elle devient
\[ \text{II.}\ \alpha' e^{x\omega} + \beta' e^{-x\omega} + \{\alpha' e^{\omega}(\text{sin.}\,\omega + \text{cos.}\,\omega) + \beta' e^{-\omega}(\text{sin.}\,\omega - \text{cos.}\,\omega)\}\,\text{sin.}\,x\omega \]
\[ + \{\alpha' e^{\omega}(\text{cos.}\,\omega - \text{sin.}\,\omega) + \beta' e^{-\omega}(\text{sin}\,\omega + \text{cos.}\,\omega)\}\,\text{cos}\,x\omega = 0 \]
ou
\[ \alpha'\{e^{x\omega} + e^{\omega}(\text{sin.}\,\omega\,\text{sin.}\,x\omega + \text{cos.}\,\omega\,\text{sin.}\,x\omega + \text{cos.}\,\omega\,\text{cos.}\,x\omega - \text{sin.}\,\omega\,\text{cos}\,x\omega)\} \]
\[ + \beta'\{e^{-x\omega} + e^{-\omega}(\text{sin.}\,\omega\,\text{sin.}\,x\omega - \text{cos.}\,\omega\,\text{sin.}\,x\omega + \text{sin.}\,\omega\,\text{cos.}\,x\omega + \text{cos.}\,\omega\,\text{cos}\,x\omega)\} = 0 \]
ou enfin
\[ \alpha'\{e^{x\omega} + e^{\omega}(\text{cos}\,(1-x)\omega - \text{sin}\,(1-x)\omega)\} + \beta'\{e^{-x\omega} + e^{-\omega}(\text{cos}\,(1-x)\omega + \text{sin}\,(1-x)\omega)\} = 0 \]
d'où on tire
\[ \beta' = -\frac{\alpha' e^{\omega}\{e^{(x-1)\omega} + \text{cos}\,(1-x)\omega - \text{sin}\,(1-x)\omega\}}{e^{-\omega}\{e^{(1-x)\omega} + \text{cos}\,(1-x)\omega + \text{sin}\,(1-x)\omega\}} \]
parconséquent
\[ \gamma' = \alpha' e^{\omega}(\text{sin.}\,\omega + \text{cos.}\,\omega) + \beta' e^{-\omega}(\text{sin.}\,\omega - \text{cos.}\,\omega) \]
\[ = \frac{\begin{array}{c}\alpha'(\text{sin.}\,\omega + \text{cos.}\,\omega)\{e^{(1-x)\omega} + \text{cos}\,(1-x)\omega + \text{sin}\,(1-x)\omega\} \\  - \alpha'(\text{sin.}\,\omega - \text{cos.}\,\omega)\{e^{(x-1)\omega} + \text{cos}\,(1-x)\omega - \text{sin}\,(1-x)\omega\}\end{array}}{e^{-\omega}\{e^{(1-x)\omega} + \text{cos}\,(1-x)\omega + \text{sin}\,(1-x)\omega\}} \]
\[ = \frac{\alpha'\{(e^{(1-x)\omega} - e^{(x-1)\omega})\,\text{sin.}\,\omega + (e^{(1-x)\omega} + e^{(x-1)\omega})\,\text{cos.}\,\omega\}}{e^{-\omega}\{e^{(1-x)\omega} + \text{cos.}\,(1-x)\omega + \text{sin}\,(1-x)\omega\}} \]
\[ + \frac{\begin{array}{c}\alpha'\{\text{sin.}\,\omega\,\text{cos.}\,(1-x)\omega + \text{sin.}\,\omega\,\text{sin.}\,(1-x)\omega + \text{cos.}\,\omega\,\text{cos}\,(1-x)\omega + \text{cos.}\,\omega\,\text{sin}\,(1-x)\omega \\  - \text{sin.}\,\omega\,\text{cos}\,(1-x)\omega + \text{sin.}\,\omega\,\text{sin}\,(1-x)\omega + \text{cos}\,\omega\,\text{cos}\,(1-x)\omega - \ill{}\}\end{array}}{e^{-\omega}\{e^{(1-x)\omega} + \text{cos}\,(1-x)\omega + \text{sin}\,(1-x)\omega\}} \]
\note{En tête, au milieu, « (9) » biffé ; en haut à droite « 164 », à
l'encre. Le paragraphe continue la solution du problème posé à la vue
334. Le numéro du paragraphe d'Euler cité à la deuxième ligne se lit
« 19 », le second chiffre à longue queue. Les équations V et VI sont
écrites sur une même ligne, VII et VIII sur la suivante. Dans la
première grande équation de $\gamma'$, l'exposant $(1-x)\omega$ du
premier crochet est surchargé. Le dernier terme de la dernière ligne,
serré contre le bord droit du feuillet, est un groupe empâté suivi de ce
qui semble « sin », et n'est pas lu ; le développement demanderait
$\text{cos}\,\omega\,\text{sin}\,(1-x)\omega$. La page s'arrête sur ce
calcul ; la suite est à la vue 338. Le verso (vue 337) est blanc.}

\page{338}

enfin
\[ \gamma' = \frac{\alpha'\{(e^{(1-x)\omega} - e^{(x-1)\omega})\,\text{sin.}\,\omega + (e^{(1-x)\omega} + e^{(x-1)\omega})\,\text{cos.}\,\omega + 2\,\text{cos.}\,x\omega\}}{e^{-\omega}\{e^{(1-x)\omega} + \text{cos}\,(1-x)\omega + \text{sin}\,(1-x)\omega\}} \]
on trouve demême
\[ \delta' = \alpha' e^{\omega}(\text{cos.}\,\omega - \text{sin.}\,\omega) + \beta' e^{-\omega}(\text{sin.}\,\omega + \text{cos.}\,\omega) \]
\[ = \frac{\begin{array}{c}\alpha'(\text{cos.}\,\omega - \text{sin.}\,\omega)\{e^{(1-x)\omega} + \text{cos}\,(1-x)\omega + \text{sin.}\,(1-x)\omega\} \\  - \alpha'(\text{sin.}\,\omega + \text{cos.}\,\omega)\{e^{(x-1)\omega} + \text{cos}\,(1-x)\omega - \text{sin}\,(1-x)\omega\}\end{array}}{e^{-\omega}\{e^{(1-x)\omega} + \text{cos}\,(1-x)\omega + \text{sin}\,(1-x)\omega\}} \]
\[ = \frac{\alpha'\{(e^{(1-x)\omega} - e^{(x-1)\omega})\,\text{cos.}\,\omega - (e^{(1-x)\omega} + e^{(x-1)\omega})\,\text{sin.}\,\omega\}}{e^{-\omega}\{e^{(1-x)\omega} + \text{cos}\,(1-x)\omega + \text{sin}\,(1-x)\omega\}} \]
\[ + \frac{\begin{array}{c}\alpha'\{(\text{cos.}\,\omega - \text{sin.}\,\omega)(\text{cos}\,(1-x)\omega + \text{sin}\,(1-x)\omega) \\  - (\text{sin}\,\omega + \text{cos}\,\omega)(\text{cos}\,(1-x)\omega - \text{sin}\,(1-x)\omega)\}\end{array}}{e^{-\omega}\{e^{(1-x)\omega} + \text{cos}\,(1-x)\omega + \text{sin}\,(1-x)\omega\}} \]
enfin
\[ \delta' = \frac{\alpha'\{(e^{(1-x)\omega} - e^{(x-1)\omega})\,\text{cos.}\,\omega - (e^{(1-x)\omega} + e^{(x-1)\omega})\,\text{sin.}\,\omega - 2\,\text{sin.}\,x\omega\}}{e^{-\omega}\{e^{(1-x)\omega} + \text{cos}\,(1-x)\omega + \text{sin.}\,(1-x)\omega\}} \]

Je prens ensuite les valeurs de $\gamma$ et $\delta$ données par les équations V et VI pour les
substituer dans l'équation I. $\alpha e^{x\omega} + \beta e^{-x\omega} + \gamma\,\text{sin}\,x\omega + \delta\,\text{cos.}\,x\omega = 0$: elle devient par là
\[ \alpha e^{x\omega} + \beta e^{-x\omega} + (\alpha - \beta)\,\text{sin.}\,x\omega + (\alpha + \beta)\,\text{cos.}\,x\omega = 0 \]
ou \ldots\ \ldots
\[ \alpha\{e^{x\omega} + \text{sin.}\,x\omega + \text{cos.}\,x\omega\} + \beta\{e^{-x\omega} - \text{sin.}\,x\omega + \text{cos.}\,x\omega\} = 0 \]
d'où on tire
\[ \beta = -\frac{\alpha\{e^{x\omega} + \text{sin.}\,x\omega + \text{cos.}\,x\omega\}}{e^{-x\omega} - \text{sin.}\,x\omega + \text{cos.}\,x\omega} \]
et parconséquent
\[ \gamma = \alpha - \beta = \frac{\alpha\{e^{-x\omega} - \text{sin.}\,x\omega + \text{cos.}\,x\omega + e^{x\omega} + \text{sin}\,x\omega + \text{cos}\,x\omega\}}{e^{-x\omega} - \text{sin.}\,x\omega + \text{cos.}\,x\omega} \]
\[ = \frac{\alpha\{e^{x\omega} + e^{-x\omega} + 2\,\text{cos.}\,x\omega\}}{e^{-x\omega} - \text{sin.}\,x\omega + \text{cos}\,x\omega} \]
\[ \delta = \alpha + \beta = \frac{\alpha\{e^{-x\omega} - \text{sin.}\,x\omega + \text{cos.}\,x\omega - e^{x\omega} - \text{sin}\,x\omega - \text{cos}\,x\omega\}}{e^{-x\omega} - \text{sin.}\,x\omega + \text{cos.}\,x\omega} \]
\[ = -\frac{\alpha\{e^{x\omega} - e^{-x\omega} + 2\,\text{sin}\,x\omega\}}{e^{-x\omega} - \text{sin}\,x\omega + \text{cos}\,x\omega} \]
Il faut àprésent substituer les valeurs que l'on vient de trouver pour les coëfficiens
$\beta, \gamma, \delta$ et $\beta', \gamma', \delta'$ dans l'équation
\[ \text{III.}\ \alpha e^{x\omega} - \beta e^{-x\omega} + \gamma\,\text{cos.}\,x\omega - \delta\,\text{sin.}\,x\omega = \alpha' e^{x\omega} - \beta' e^{-x\omega} + \gamma'\,\text{cos.}\,x\omega - \delta'\,\text{sin.}\,x\omega. \]
\note{En tête, au milieu, « (10) » biffé ; en haut à droite « 165 », à
l'encre. La page reprend le calcul de $\gamma'$ laissé au bas de la vue
336. Le « sin » du premier crochet de la deuxième expression de
$\delta'$ est surchargé, comme « pour » à l'avant-dernière ligne. Après
« ou », à la ligne qui suit l'équation I transformée, une ligne de points
mène à l'équation ; elle est rendue par des points de suspension. La
page s'arrête sur l'équation III, sans conclure ; la suite n'est pas
dans ce lot (voir la vue 340). Le verso (vue 339) est blanc.}

\page{340}

par cette substitution le premier membre devient:
\[ \frac{\begin{array}{c}\alpha\{1 - e^{x\omega}(\text{sin.}\,x\omega - \text{cos.}\,x\omega)\} + \alpha\{1 + e^{-x\omega}(\text{sin.}\,x\omega + \text{cos.}\,x\omega)\} \\  + \alpha\{e^{x\omega} + e^{-x\omega} + 2\,\text{cos}\,x\omega\}\,\text{cos}\,x\omega + \alpha\{e^{x\omega} - e^{-x\omega} + 2\,\text{sin.}\,x\omega\}\,\text{sin.}\,x\omega\end{array}}{e^{-x\omega} - \text{sin.}\,x\omega + \text{cos.}\,x\omega} \]
\[ = \frac{\begin{array}{c}\alpha\{4 - (e^{x\omega} - e^{-x\omega})\,\text{sin.}\,x\omega + (e^{x\omega} + e^{-x\omega})\,\text{cos.}\,x\omega \\  + (e^{x\omega} + e^{-x\omega})\,\text{cos.}\,x\omega + (e^{x\omega} - e^{-x\omega})\,\text{sin.}\,x\omega\}\end{array}}{e^{-x\omega} - \text{sin.}\,x\omega + \text{cos.}\,x\omega} \]
\[ = \frac{2\alpha\{2 + (e^{x\omega} + e^{-x\omega})\,\text{cos}\,x\omega\}}{e^{-x\omega} - \text{sin}\,x\omega + \text{cos.}\,x\omega} \]
et le second membre de la même équation
\[ \frac{\begin{array}{c}\alpha' e^{(x-1)\omega}\{e^{(1-x)\omega} + \text{cos}\,(1-x)\omega + \text{sin}\,(1-x)\omega\} \\  + \alpha' e^{(1-x)\omega}\{e^{(x-1)\omega} + \text{cos}\,(1-x)\omega - \text{sin}\,(1-x)\omega\}\end{array}}{e^{-\omega}\{e^{(1-x)\omega} + \text{cos}\,(1-x)\omega + \text{sin}\,(1-x)\omega\}} \]
\[ + \frac{\alpha'\{(e^{(1-x)\omega} - e^{(x-1)\omega})\,\text{sin.}\,\omega + (e^{(1-x)\omega} + e^{(x-1)\omega})\,\text{cos.}\,\omega + 2\,\text{cos}\,x\omega\}\,\text{cos.}\,x\omega}{e^{-\omega}\{e^{(1-x)\omega} + \text{cos}\,(1-x)\omega + \text{sin}\,(1-x)\omega\}} \]
\[ - \frac{\alpha'\{(e^{(1-x)\omega} - e^{(x-1)\omega})\,\text{cos.}\,\omega - (e^{(1-x)\omega} + e^{(x-1)\omega})\,\text{sin.}\,\omega - 2\,\text{sin.}\,x\omega\}\,\text{sin.}\,x\omega}{e^{-\omega}\{e^{(1-x)\omega} + \text{cos}\,(1-x)\omega + \text{sin}\,(1-x)\omega\}} \]
\[ = \frac{\alpha'\{4 + (e^{(1-x)\omega} + e^{(x-1)\omega})\,\text{cos}\,(1-x)\omega - (e^{(1-x)\omega} - e^{(x-1)\omega})\,\text{sin}\,(1-x)\omega\}}{e^{-\omega}\{e^{(1-x)\omega} + \text{cos}\,(1-x)\omega + \text{sin}\,(1-x)\omega\}} \]
\[ + \frac{\begin{array}{c}\alpha'\{(e^{(1-x)\omega} + e^{(x-1)\omega})[\text{cos.}\,x\omega\,\text{cos.}\,\omega + \text{sin.}\,x\omega\,\text{sin.}\,\omega) \\  + (e^{(1-x)\omega} - e^{(x-1)\omega})(\text{sin.}\,\omega\,\text{cos.}\,x\omega - \text{cos.}\,\omega\,\text{sin.}\,x\omega)\}\end{array}}{e^{-\omega}\{e^{(1-x)\omega} + \text{cos}\,(1-x)\omega + \text{sin}\,(1-x)\omega\}} \]
\[ = \frac{2\alpha'\{2 + (e^{(1-x)\omega} + e^{(x-1)\omega})\,\text{cos}\,(1-x)\omega\}}{e^{-\omega}\{e^{(1-x)\omega} + \text{cos}\,(1-x)\omega + \text{sin}\,(1-x)\omega\}} \]
Ainsi l'équation III se réduit, en dernière analyse, à
\[ \frac{\alpha\{2 + (e^{x\omega} + e^{-x\omega})\,\text{cos}\,x\omega\}}{e^{-x\omega} - \text{sin.}\,x\omega + \text{cos.}\,x\omega} = \frac{\alpha'\{2 + (e^{(1-x)\omega} + e^{(x-1)\omega})\,\text{cos}\,(1-x)\omega\}}{e^{-\omega}\{e^{(1-x)\omega} + \text{cos}\,(1-x)\omega + \text{sin}\,(1-x)\omega\}} \]
ou à
\[ \frac{\alpha e^{-\omega}\{e^{(1-x)\omega} + \text{cos}\,(1-x)\omega + \text{sin}\,(1-x)\omega\}}{\alpha'\{e^{-x\omega} - \text{sin.}\,x\omega + \text{cos.}\,x\omega\}} = \frac{2 + (e^{(1-x)\omega} + e^{(x-1)\omega})\,\text{cos}\,(1-x)\omega}{2 + (e^{x\omega} + e^{-x\omega})\,\text{cos.}\,x\omega} \]
Il faut encore avoir recours à l'équation IV, que l'on peut écrire comme il suit:
\[ \alpha e^{x\omega} + \beta e^{-x\omega} - \gamma\,\text{sin.}\,x\omega - \delta\,\text{cos.}\,x\omega = \alpha' e^{x\omega} + \beta' e^{-x\omega} - \gamma'\,\text{sin.}\,x\omega - \delta'\,\text{cos}\,x\omega \]
\note{En tête, au milieu, « (11) » biffé ; en haut à droite « 166 », à
l'encre. La page poursuit la substitution annoncée au bas de la vue 338.
Le « 4 » par lequel s'ouvrent les deux accolades « $\alpha\{4 - \ldots$ »
et « $\alpha'\{4 + \ldots$ » est tracé sur un autre chiffre ; le « sin »
du dénominateur de la première fraction est surchargé. Le crochet ouvert
« [ » devant « $\text{cos.}\,x\omega\,\text{cos.}\,\omega$ » est fermé par
une parenthèse ; la leçon est gardée. « l'équation IV » : le I est un
trait court collé au V. Le feuillet s'arrête sur l'équation IV réécrite,
au bas de la page ; la suite n'est pas dans ce lot. Le verso (vue 341)
est blanc, avec le texte de la vue 340 visible par transparence.}

\end{document}
